titulo: Teorema de Heine-Borel
tema: Compacidad
dificultad: avanzada
etiquetas: compacidad, espacios métricos, cubrimientos
actualizada: 2026-07-30
---
\begin{teorema}
Un subconjunto $K \subseteq \mathbb{R}^n$ es compacto si y sólo si es cerrado y acotado.
\end{teorema}

\begin{demostracion}
Esbozo. Si $K$ es compacto, el cubrimiento por bolas $B(0,m)$ con $m \in \mathbb{N}$
admite subcubrimiento finito, luego $K$ es acotado; y el complemento de $K$ es
abierto separando puntos, luego $K$ es cerrado.

Recíprocamente, si $K$ es cerrado y acotado está contenido en una celda
$[-r,r]^n$, que es compacta por bisección sucesiva, y todo cerrado dentro de un
compacto es compacto. \qed
\end{demostracion}
